\section{Fixed bounded-formula reflection}
\label{sec:bounded-reflection}

Bounded formulas are the main bridge between ordinary finite reasoning and quantifier-free induction.  The compiler described here is a schema for each fixed formula.  It is not an object-level interpreter for arbitrary syntax codes.

\subsection{The bounded syntax}

Fix free variables \(\vec x\).  Let \(\mathcal B\) be generated by:
\[
\begin{aligned}
A ::= {}& s(\vec x)=t(\vec x) \mid \bot \mid A\wedge A \mid A\vee A
\mid A\to A \\
&\mid \exists y<B(\vec x)\,A(y,\vec x)
\mid \forall y<B(\vec x)\,A(y,\vec x),
\end{aligned}
\]
where every term denotes a fixed primitive-recursive function.  The bound is independent of the newly bound variable.  In the actual first-order syntax, the two bounded quantifiers abbreviate
\[
\exists y\,(y<B\wedge A),
\qquad
\forall y\,(y<B\to A).
\]

Implication is included because application theorems naturally use it.  Its Boolean compilation is justified constructively because every formula in \(\mathcal B\) will be proved decidable.

\subsection{Boolean terms}

A term \(b\) is a \emph{proved bit} when iPRA derives \(b=0\vee b=1\).  For proved bits define
\[
\begin{aligned}
\neg_b b &= 1\mathbin{\dot-}b,\\
b\wedge_b c &= bc,\\
b\vee_b c &= 1\mathbin{\dot-}\bigl((1\mathbin{\dot-}b)(1\mathbin{\dot-}c)\bigr),\\
b\Rightarrow_b c &= (\neg_b b)\vee_b c.
\end{aligned}
\]
Here \(\dot-\) is truncated subtraction.  All operations are primitive recursive.

\begin{lemma}[Boolean truth tables]
For proved bits \(b,c\), iPRA derives that the displayed operations are bits and satisfy the usual truth tables.  In particular,
\[
\begin{aligned}
 b\wedge_b c=1 &\leftrightarrow b=1\wedge c=1,\\
 b\vee_b c=1 &\leftrightarrow b=1\vee c=1,\\
 \neg_b b=1 &\leftrightarrow b=0,\\
 b\Rightarrow_b c=1 &\leftrightarrow (b=0\vee c=1).
\end{aligned}
\]
\end{lemma}
\begin{proof}
Use the two supplied disjunctions \(b=0\vee b=1\) and \(c=0\vee c=1\).  There are four finite cases.  In each case the result follows from the checked defining equations and semiring/monus laws.  This is intuitionistic case elimination on proved finite alternatives, not unrestricted excluded middle.
\end{proof}

Atomic equality is compiled by the checked term \(\mathsf{eqb}(s,t)\).  The current library proves \(s=t\to\mathsf{eqb}(s,t)=1\) and \(s\ne t\to\mathsf{eqb}(s,t)=0\).  Combined with the checked decidability of term equality, this gives the atomic reflection package.

\subsection{Arbitrary-arity finite folds}

Let \(b(i,\vec x)\) be a fixed bit term.  Define
\[
\begin{aligned}
\mathsf{All}_b(0,\vec x)&=1,&
\mathsf{All}_b(n+1,\vec x)&=\mathsf{All}_b(n,\vec x)b(n,\vec x),\\
\mathsf{Any}_b(0,\vec x)&=0,&
\mathsf{Any}_b(n+1,\vec x)&=\mathsf{Any}_b(n,\vec x)\vee_b b(n,\vec x).
\end{aligned}
\]

These definitions work with any fixed number of parameters.  The implementation uses \code{PTm.toPR} to abstract the open body and \code{recDef} to introduce the recursion argument and accumulator.  The parameter-renaming map sends the body's index variable to step variable \(0\), skips the accumulator at step variable \(1\), and sends every outer parameter two positions upward.  This construction is written explicitly in \code{lean/SCOPE/BoundedFormulaBlueprint.lean}.

\begin{lemma}[Finite existential reflection]
\label{lem:any-reflection}
For a proved bit body,
\[
  \mathsf{Any}_b(n,\vec x)=1
  \quad\leftrightarrow\quad
  \exists i<n\; b(i,\vec x)=1.
\]
Moreover a least such \(i\) is primitive recursive, with \(n\) as the failure sentinel.
\end{lemma}
\begin{proof}
Search for a zero of the complement \(1\mathbin{\dot-}b(i,\vec x)\).  The checked search term \(\mu(n,\vec x)\) satisfies \(\mu\le n\); if \(\mu<n\), then the complement at \(\mu\) is zero, hence \(b(\mu)=1\); and every earlier index has nonzero complement, hence bit value zero for the original body is impossible there.  If a witness \(i<n\) with \(b(i)=1\) is supplied, \code{search_found} yields \(\mu<n\), validity at \(\mu\), and \(\mu\le i\).

It remains to connect the recursive disjunction fold to the search success bit.  Prove by quantifier-free induction on \(n\) the equation
\[
  \mathsf{Any}_b(n)=1
  \quad\leftrightarrow\quad
  \mu(n)<n.
\]
The induction matrix can be expressed as an equality of characteristic terms.  At the successor step, either a witness was already found, in which case both sides remain true, or the sentinel was \(n\); then both sides become true exactly when \(b(n)=1\).  All branching is on reflected term comparisons and body bits.  Combining the equation with the checked search specification gives the desired existential equivalence and leastness.
\end{proof}

\begin{lemma}[Finite universal reflection]
\label{lem:all-reflection}
For a proved bit body,
\[
  \mathsf{All}_b(n,\vec x)=1
  \quad\leftrightarrow\quad
  \forall i<n\;b(i,\vec x)=1.
\]
\end{lemma}
\begin{proof}
Search for a counterexample, that is, an index with \(b(i)=0\).  Equivalently search for a zero of \(b\) itself.  Let \(\nu(n)\) be the least zero or sentinel \(n\).  A quantifier-free induction on the recursion equations proves
\[
  \mathsf{All}_b(n)=1 \quad\leftrightarrow\quad \nu(n)=n.
\]

For the forward implication, suppose \(i<n\).  If \(b(i)=0\), checked search foundness gives \(\nu<n\), contradicting \(\nu=n\).  Since \(b(i)\) is a bit, it must equal one.

For the backward implication, assume every \(i<n\) has value one.  If \(\nu<n\), the search specification gives \(b(\nu)=0\), while the universal premise instantiated at \(\nu\) gives \(b(\nu)=1\), a contradiction.  Therefore \(\nu=n\), and the fold equation yields \(\mathsf{All}_b(n)=1\).

No induction predicate contains the bounded universal formula.  The only induction is on equations defining the fold and search state.
\end{proof}

\subsection{The reflection theorem}

\begin{theorem}[Fixed bounded reflection]
\label{thm:bounded-reflection}
For every fixed \(A(\vec x)\in\mathcal B\), there is a primitive-recursive term \(\bit_A(\vec x)\) and iPRA derivations, uniform in the free parameters, of
\[
\begin{aligned}
&\bit_A=0\vee\bit_A=1,\\
&A\to\bit_A=1,\\
&\bit_A=1\to A.
\end{aligned}
\]
Consequently iPRA proves \(A\leftrightarrow\bit_A=1\) and \(A\vee\neg A\).
\end{theorem}

\begin{proof}
Proceed by structural recursion on the fixed syntax tree.

For an equality \(s=t\), take \(\bit_A=\mathsf{eqb}(s,t)\).  Equality decidability and the two checked \code{eqb} lemmas provide Booleanity and both implications.  For \(\bot\), take the constant zero bit.  The forward implication follows by bottom elimination and the backward implication is \(0=1\to\bot\), using successor separation.

For conjunction and disjunction, apply the induction hypotheses and the Boolean truth-table lemma.  For example, from \(A\wedge B\) obtain both component bits equal to one and hence their product equal to one.  Conversely, product one implies both bits one; reflect them back and introduce the conjunction.

For implication, let \(a=\bit_A\) and \(b=\bit_B\).  Suppose \(A\to B\).  Decide \(A\) using the induction hypothesis.  If \(\neg A\), then \(a=0\) and the implication bit is one.  If \(A\), obtain \(B\), hence \(b=1\), and again the implication bit is one.  Conversely, if \(a\Rightarrow_b b=1\), the truth table gives \(a=0\vee b=1\).  Under an assumption \(A\), the first alternative contradicts the forward reflection of \(A\); the second reflects to \(B\).  Thus \(A\to B\) follows intuitionistically.

For \(\exists y<B\,A(y)\), compile the body to \(b(y)\) and take \(\bit=\mathsf{Any}_b(B)\).  A supplied existential witness gives \(i<B\) and \(A(i)\), hence \(b(i)=1\), so \cref{lem:any-reflection} gives the aggregate bit.  Conversely, aggregate bit one supplies the least \(i<B\) with \(b(i)=1\); the induction hypothesis yields \(A(i)\), and existential introduction finishes.

For \(\forall y<B\,A(y)\), take \(\mathsf{All}_b(B)\).  A universal proof gives each body bit one, hence the fold bit by \cref{lem:all-reflection}.  Conversely, fold bit one gives \(b(i)=1\) for an arbitrary \(i<B\); reflect to \(A(i)\) and universally introduce the bound variable.

Booleanity is preserved at every constructor.  Finally, decide whether \(\bit_A=1\).  In the positive case reflect to \(A\).  In the negative case, Booleanity forces \(\bit_A=0\); if \(A\) held, forward reflection would give \(\bit_A=1\), contradiction.  Therefore \(A\vee\neg A\).
\end{proof}

\subsection{Binder-safe specialization}

The theorem is stated over open formulas, so substitution of arbitrary terms is obtained directly from \code{Prf.subst}.  Beneath a quantifier, \code{Syn.up} preserves the newly bound variable and shifts the substituted outer terms.  The checked identities \code{Fm.shift_subst_up} and \code{Fm.subst_inst} prove that weakening and instantiation commute in the required form.

Thus the first implementation milestone does not need a bespoke semantic substitution proof for every application.  It needs a syntactic compiler coherence theorem only to show that regenerated code is extensionally equal to the substituted old code.  The already checked \code{toPR_eq} theorem supplies the key step for terms abstracted into primitive-recursive symbols.

\subsection{Canonical guarded bounded choice}

\begin{theorem}[Least accepted witness]
\label{thm:least-accepted}
Let \(R(\vec x,y)\in\mathcal B\), and let \(B(\vec x)\) be a primitive-recursive bound.  There are primitive-recursive terms \(w(\vec x)\) and \(h(\vec x)\in\{0,1\}\) such that iPRA proves
\[
\begin{aligned}
 h=1 &\leftrightarrow \exists y<B\;R(\vec x,y),\\
 h=1 &\to w<B\wedge R(\vec x,w),\\
 h=0 &\to w=B,\\
 h=1\wedge z<w &\to \neg R(\vec x,z).
\end{aligned}
\]
\end{theorem}
\begin{proof}
Compile \(R\) to its bit \(r(y,\vec x)\).  Apply the checked least-zero search to \(1\mathbin{\dot-}r\).  Let \(w\) be the search result and let \(h\) be the reflected comparison bit for \(w<B\).  The search specification gives the sentinel law, soundness, and minimality.  The reflection theorem translates between \(r=1\) and \(R\).  Every displayed claim follows by finite intuitionistic reasoning.
\end{proof}

\begin{corollary}[Bounded choice schema]
If iPRA proves \(\forall\vec x\,\exists y<B(\vec x)\,R(\vec x,y)\), then it proves \(\forall\vec x\,R(\vec x,w(\vec x))\) for the concrete primitive-recursive least witness term \(w\).
\end{corollary}

The bound is essential.  This is not countable choice and does not turn arbitrary total relations into primitive-recursive functions.

\subsection{Why bounded reflection is enough for many classical finite arguments}

A finite decidable argument may be discovered classically, by SAT solving, computer algebra, or exhaustive search.  To import it constructively, expose a finite certificate and a bounded verifier.  The certificate need not be canonical unless downstream syntax needs its code.  In the common case, private sealing eliminates it after the verifier derives the certificate-independent theorem.

This separation permits aggressive external discovery while retaining a small trusted object theory: only the verifier equations and their \(\Prf\,\PRAAx\) correctness proofs matter to the final theorem.
